\begin{proofbox}
\textbf{统一记号。}设 $P(x_0,y_0)$，$F_1(-c,0)$、$F_2(c,0)$，
$r_1=PF_1$、$r_2=PF_2$，$\theta=\angle F_1PF_2$，
$e=c/a$。由椭圆定义及焦半径公式，
\[
 r_1+r_2=2a,\qquad r_1=a+ex_0,\qquad r_2=a-ex_0,
 \qquad c^2=a^2-b^2.
\]
先在 $y_0\ne0$ 时把 $F_1PF_2$ 视为非退化三角形；
$y_0=0$ 的两个长轴端点将在最后单独验证。

\medskip
\textbf{方法一：余弦定理 $\boldsymbol{\to}$ 半角公式（代数通法）。}
\begin{enumerate}
  \item 在 $\triangle F_1PF_2$ 中，顶角 $\theta$ 的对边为 $2c$。
  由余弦定理
  \[
  (2c)^2=r_1^2+r_2^2-2r_1r_2\cos\theta,
  \quad
  \cos\theta=\frac{r_1^2+r_2^2-4c^2}{2r_1r_2}.
  \]

  \item 代入 $r_1=a+ex_0$、$r_2=a-ex_0$，得到
  \[
  r_1r_2=a^2-e^2x_0^2,\qquad
  r_1^2+r_2^2=2(a^2+e^2x_0^2),
  \]
  从而
  \[
  \cos\theta
  =1-\frac{2(c^2-e^2x_0^2)}{a^2-e^2x_0^2}.
  \]

  \item 由椭圆方程 $x_0^2=a^2(1-y_0^2/b^2)$ 及 $e=c/a$，
  \[
  c^2-e^2x_0^2=\frac{c^2y_0^2}{b^2},\qquad
  a^2-e^2x_0^2=\frac{b^4+c^2y_0^2}{b^2}.
  \]
  因此
  \[
  \boxed{\cos\theta=\frac{b^4-c^2y_0^2}{b^4+c^2y_0^2}}.
  \]

  \item 采用半角恒等式（$0\leq\theta<180^\circ$）
  \[
  \tan^2\frac\theta2=\frac{1-\cos\theta}{1+\cos\theta}
  =\frac{c^2y_0^2}{b^4}.
  \]
  由于 $\tan(\theta/2)\geq0$，故
  \[
  \boxed{\tan\frac\theta2=\frac{c|y_0|}{b^2}}.
  \]
\end{enumerate}

\medskip
\textbf{方法二：面积 + 内切圆（几何直观法）。}
\begin{enumerate}
  \item 令焦点三角形面积为 $S$、半周长为 $s$、内切圆半径为 $\rho$。
  由于底边 $F_1F_2=2c$，高为 $|y_0|$，
  \[
   S=c|y_0|,\qquad
   s=\frac{r_1+r_2+2c}{2}=a+c,\qquad
   \rho=\frac S s=\frac{c|y_0|}{a+c}.
  \]

  \item 设内切圆与边 $PF_1$ 切于 $T$，内心为 $I$。
  由同一点引圆的切线段相等，
  \[
  PT=s-F_1F_2=(a+c)-2c=a-c.
  \]
  $PI$ 平分顶角，$IT\perp PF_1$，因此在直角三角形 $PIT$ 中
  \[
  \tan\frac\theta2=\frac{IT}{PT}
  =\frac{\rho}{a-c}
  =\frac{c|y_0|}{(a+c)(a-c)}
  =\boxed{\frac{c|y_0|}{b^2}}.
  \]
  \item \textbf{本质：}焦点三角形的底边、半周长和切线段 $PT=a-c$
  都固定；增大 $|y_0|$ 就增大面积 $S$、内切圆半径 $\rho$，
  从而增大 $\tan(\theta/2)$ 和顶角 $\theta$。
\end{enumerate}

\medskip
\textbf{方法三：正弦面积式 + 余弦定理（交叉核验）。}
\begin{enumerate}
  \item 对同一焦点三角形，用两边及夹角求面积：
  \[
  S=\frac12 r_1r_2\sin\theta=c|y_0|,
  \quad\Longrightarrow\quad
  \sin\theta=\frac{2c|y_0|}{r_1r_2}.
  \]
  \item 另一方面，余弦定理及 $r_1+r_2=2a$ 给出
  \[
  1+\cos\theta
  =\frac{(r_1+r_2)^2-(2c)^2}{2r_1r_2}
  =\frac{4(a^2-c^2)}{2r_1r_2}
  =\frac{2b^2}{r_1r_2}.
  \]
  \item 使用 $\tan(\theta/2)=\sin\theta/(1+\cos\theta)$，得
  \[
  \boxed{\tan\frac\theta2
  =\frac{2c|y_0|/(r_1r_2)}{2b^2/(r_1r_2)}
  =\frac{c|y_0|}{b^2}}.
  \]
  此法不需要解出 $x_0$，适合直接把三角形面积、正弦与余弦定理串联。
\end{enumerate}

\medskip
\textbf{补充推论：由余弦定理直接推出面积公式（固定角决定面积）。}
这里再换一个解题视角：\textbf{先固定顶角 $\theta$，再考察焦点三角形的面积。}
由余弦定理和椭圆定义，
\begin{align*}
  4c^2
  &=r_1^2+r_2^2-2r_1r_2\cos\theta\\
  &=(r_1+r_2)^2-2r_1r_2(1+\cos\theta)\\
  &=4a^2-4r_1r_2\cos^2\frac\theta2.
\end{align*}
整理得（因 $0\leq\theta<180^\circ$，半角余弦不为零）
\[
  \boxed{r_1r_2=\frac{a^2-c^2}{\cos^2(\theta/2)}
  =\frac{b^2}{\cos^2(\theta/2)}}.
\]
于是由两边及夹角的面积公式，
\begin{align*}
  S
  &=\frac12r_1r_2\sin\theta\\
  &=\frac12\cdot\frac{b^2}{\cos^2(\theta/2)}
  \cdot2\sin\frac\theta2\cos\frac\theta2\\
  &=\boxed{b^2\tan\frac\theta2}.
\end{align*}
又 $S=c|y_0|$，故
\[
  \boxed{\theta\ \text{相同}\iff S\ \text{相同}
  \iff |y_0|\ \text{相同}}.
\]
这一等价关系成立的前提是\textbf{同一个固定的椭圆}（$a,b,c$
不改变），且两腰和 $r_1+r_2=2a$ 始终不变。

\medskip
\textbf{定角反求点：一般有四个对称点。}
给定 $0\leq\theta\leq\theta_{\max}$，由上式先得
\[
  |y_0|=\frac{b^2}{c}\tan\frac\theta2,
  \qquad
  |x_0|=a\sqrt{1-\frac{y_0^2}{b^2}}.
\]
当 $0<\theta<\theta_{\max}$ 时，两坐标绝对值都大于零，
点 $P=(\pm|x_0|,\pm|y_0|)$ 通常有四个（两个正负号独立选取）；
当 $\theta=0$ 时只有 $P=(\pm a,0)$；
当 $\theta=\theta_{\max}$ 时只有 $P=(0,\pm b)$。
这也严格说明了\textbf{焦角决定坐标绝对值，却不唯一决定点的位置}。

\medskip
\textbf{最大角、最小角的确定与端点检验。}
由椭圆方程 $|y_0|\leq b$，等号当且仅当 $P=(0,\pm b)$。
又因 $\theta/2\in[0,90^\circ)$，$\tan$ 严格递增，故
\[
 \boxed{\theta_{\max}=2\arctan\frac cb}.
\]
在 $P=(0,b)$ 时，$PF_1=PF_2=a$，$PO=b$、$OF_2=c$，
由等腰三角形顶角平分线所形成的直角三角形直接得
\[
 \tan\frac{\theta_{\max}}2=\frac cb,\quad
 \sin\frac{\theta_{\max}}2=\frac ca=e,\quad
 \cos\theta_{\max}=1-2e^2=\frac{b^2-c^2}{b^2+c^2}.
\]
当 $P=(\pm a,0)$ 时，$P,F_1,F_2$ 共线且从 $P$ 发出的两条射线
方向相同，所以 $\theta=0$，与半角公式（右端为 $0$）一致。
反之，$\theta=0$ 则 $c|y_0|/b^2=0$，由 $c>0$ 得 $y_0=0$，
椭圆上只能是 $P=(\pm a,0)$。所以最小焦角为
\[
  \boxed{\theta_{\min}=0^\circ\quad\Longleftrightarrow\quad
  P=(\pm a,0)}.
\]
\textbf{因此最大值、最小值与半角结论对椭圆上所有点均成立。}
\end{proofbox}
